\chapter{Introduction} \label{ch:intro}
\index{Introduction@\emph{Introduction}}

\section{Why this book?}\label{sec:whyThisBook}
I've worked in tech field for 10 years.
But only recently it came to my realization that high-end software developers have remained rare
commodities all along.
This is corroborated by a quick peek at repositories at \url{github.com}, a search of open-source
projects that are bread-n-butters for your daily job such as Hadoop, Spark, Apache Common
Libraries, \textit{etc.}, code review of colleagues working at top-tier tech companies, or
during interview evaluations.
Wherever you look, it's not hard to spot the various `smells' as listed in Martin Fowler's 1999
seminal book~\cite{Fowler1999}.
As a matter of fact, many projects that are depended by millions of tech practitioners and myriad of
systems came with a great design but implemented with mediocrity.
These plain facts convey the crying signal that there is a lack of supply of quality engineers.
You are right, this is my personal testimony.
According to a research by Forrester, Rather, the gap is larger each year in the past 20 years.
In Forrester's 2018 report, it cites that software engineering talent shortage is quality, not
just quantity.
Forrester also projects that firms will pay $20\%$ above market for quality engineering talent in
2018 \footnote{\url{https://www.lingolive.com/2018s-software-engineering-talent-shortage}}

As pointed out in the report, the current shortage is not a lack of ``engineers'', but
one of quality-a lack of experienced engineers with a deep understanding of software engineering.

There are beliefs flowing around that emphasizes top-level designs than low-level implementation;
that architecture design, interface designs take precedence;
that implementations can be put off and delegated to less-qualified engineers.
But I would argue that
those beliefs that emphasize design over implementation are superficial observations rather than
insights.
Implementations are the key deciding factor for the success of a project.
They are just equally, if not more, important as design.
%It's the lack of quality engineers that cause this division of works.
It is more a do-or-die makeshift than willful choice-There are just not enough top-tier engineers
working on the projects.
The few top-notch minds focused their energy on something else.
Ever since software industry became an industry, the demand for qualified professionals far
outgrows the supply.

I've worked with graduates and post-graduates from top-tier schools.
They all come with latest theories in computer sciences.
Their resumes are full of eye-catching background and GPA's.
But at the end of the day when ideas need to be converted to code, the incompetence in them
manifests itself brightly.

Formal computer science education is a foundation but not the commodity of trade.
The application of education to problem solving is a skill that is learned by experience and
mentorship.
It takes years of real-world combat training to achieve this level of proficiency.
This book would not and is never my pursuance to help you cut short the path.

What I want to achieve through this book is that one may be inspired to immerse oneself in a
learning mindset from one's own coding exercise.
Through this book, I hope to do my part in narrowing the gap between demand and supply of quality
talents.
I rightly believe readers share the aspiration as mine and are working hard to realize their
potential to become better programming talents.
Also I encourage readers to form the habit of summing up their work regularly and share with
each other their learning often.

\section{Who should read this book?} \label{sec:whoshouldread}
This book is intended for professionals who already have a solid command on
algorithms, data structures, operating systems, and programming languages of their choice.
However applying these to a complex coding practice on a daily basis remains a challenging task.
toward their software engineering career goal.

This book is rather a summary of 10 years software engineering experience of the author.
The author believes and have proved through personal experience the effectiveness of constant
learning.
Undeniably, a lot of activities help one improve on programming skills.
Coding up a great deal of problems will get you familiar with various algorithms, a command of
data structures, and brush up on your coding skills.
However, an often ignored activity is to contemplate on algorithms and your implementations of them.
From time to time, one need to look back on the works they did before and pause on
anything that may catch one's eyes.
Trying new things that ranges from mundane variable renaming to systematic code overhaul.

But the beauty of this kind of activities lies in the gifts that come with practice.
The more you explore, the more there is to discover.
If you make this a habit, you will unearth many gems in the depth of your ability.
You will be a perpetual generator of innovations, a powerhouse of self-sustained improvement.

Below is just a list of things that I usually spend time for meditation:
\begin{itemize}
    \item Bugs made previously
    \item Bugs fixed previously
    \item Why a particular implementation works
    \item Why a particular implementation does not work
    \item Refactor redundant code
    \item
\end{itemize}

\section{Who am I?}\label{sec:whoAmI}
Honestly, I'm nobody.
I wish someone else could have authored this book instead of me.
But I waited for a long time, searched the corners of libraries and bookstores.
I just don't seem to see a book with similar coverage.

If you demand, I admit that I am a self-taught software engineer, converted from a little
bit of scientific research background.
You guessed right!
I didn't achieve much in research, either.
So that little research background doesn't lend me much credentials for what I~am doing here.

However I don't feel like I need the borrowed credential to write this book.
I write this book because I feel I have something to share with you.
In fact, I've got get something amusing for you to read on.
I struggled for 10 years, from the very beginning of my career as programmer to get to the bottom
of everything I wanted to while in practice.
My immersion in this field was materialized in thousands of pages of notes and dairies, $40k$
lines of solutions for coding problems wherever I can find\footnote{\textit{e.g.},
\url{https://projecteuler.net}, \url{https://leetcode.com}} (many of problems were revisited over
a period of several years for four to five times), and $100k+$ lines of deployed production code
at top-tier tech firms.
I believe my learning is worth sharing.

New book thoughts quotes.
I like what Steve McConnell in his book code complete ``the primary subject of communication is
not computer but other humans''.
My primary goal of design and coding is not bug free, but to have
fewer places for bug to hide.
This book is a zoomed-in view of the chapter in code complete.

The hands-on experience shared here will not eliminate bugs from your code.
But they will help eliminate hiding places for bugs.

\section{Hierarchical Software System}\label{sec:hierarchicalSoftwareSystem}
A fact about majority modern software system consists a tall hierarchy of many components-from
screen displays down to silicon substrates.
Many software programmer even not realized is that our understanding of what happens beneath the
execution of your familiar code is quite shallow.
If one peeks under the hood, he will come to realize that the little cute function one implemented
actually invokes a great deal of open-source libraries, frameworks, compilers, interpreters,
virtual machines, operating systems, network protocols, ROM-bore firmwares, etc.
The a simple transaction you initiate were handled by a myriad of components among the
hierarchy in a coordinated way.
With this revelation, one should be in a better position to write his own functions, classes,
APIs, and be fully motivated to make their as modular, reusable as possible.
Your functions should be built upon one another.
Together they tell a complex and complete story, with outlines as well as details at the right place.
Our world is run on a conglomeration of hierarchical, systems built on top one another, made of
various code of all varieties, large or small, open source and proprietary.
Bugs can hide at all the levels in the hierarchy.
It's truly a miracle that our software world runs at all this very moment.

The idea of remove all the bugs is ridiculously unrealistic.
Because some of the bug is burnt into the chips you are using.
Since electrons follow principle of quantum mechanics, uncertainty principles
guarantees you that your programs works only probabilistically.
That is to say, some bugs are intrinsic to the computing hardware we are using today.

This book only outlines heuristics that inspires like minds with desire for continuous improvement.
It's not a recipe to craft bug free code.
A machine wouldn't be able to learn from my book to produce better quality code in the
foreseeable future.
But it outlines a formula for like minded programmers to strive for ever better-quality code.
It is a recipe for self-improvement.

\section{What's last mile in software development}\label{sec:whatsLastMileInSoftwareDevelopment}
Design is largely done.
I'm not saying that the design is water-tight and written in a concrete slab.
It's about implementing functions.

\section{Structure}\label{sec:structure}

Why do I write this book?

From time to time, I will use interview problems to illustrate the strategy of implementation of
certain algorithms, but explaining coding questions or their solutions is never the focus of this
book.
The strategy of implementing functions and classes is applicable to both coding questions
during an interview and in everyday working environment.

Many gurus from academia as well as top-tier industry has classical works regarding programming
and software development best practice.
I assume readers are already familiar with that stuff.
What I am sharing instead is first-hand personal experience as a programmer in the past decade.
They shouldn't be viewed as rules to be accepted but to be felt and verified in your own practice.
Please provide me your feedback as you read and I'd love to hear from you.

\section{Structure of the Book}\label{sec:structureOfTheBook}
The book is divided into three parts.

In Part~\ref{prt:flowControl} ``\nameref{prt:flowControl}'', we go in length talking about the
foundations for concise and readable flow controls.
We first talk about well-structured, modular, and reusable functions.
We then dash to the meat of the essence of readable and succinct flow control technologies.

In Part~\ref{prt:dataStructure} ``\nameref{prt:dataStructure}'', we will explain the usage
patterns of common data structures such as maps, linked lists, binary trees, stacks, and queues.
For each of these data structures, I assume you already have a solid understanding of what they
are.
With a plenty of examples, I will explain how I come to realize the many usage patterns that can
help you write reliable, easily maintainable code.

With the foundations laid, We will touch base on a few advanced topics in Part~\ref{prt:blastOff}
``\nameref{prt:blastOff}''.
We will start with how to simplify embedded loops by introducing a term ``loop thinning''.
Next we will talk about sentinels and their application to reduce redundant and / or
duplicate code.
Another topic I'd like to elaborate on is dynamic programming.
I will show you how the right mindset can help you psychologically.
I will also explain the two common approaches to crack a complex dynamic programming problem.
Furthermore, I go into details to explain a multi-variate dynamic programming techniques for
solving complicated dynamic programming problems.
Yet another important topic I'd like to eloberate is object-oriented design (OOD).
When a conventional book talks about OOD, it is often done in a detached way.
Seldom anyone talks about the interaction between OOD and implementation.
In a section, I will attempt to explain the interplay between implementation and low-cost,
low-level OOD optimization.
Finally, I will dedicate a section to explain a variational method to gradually approach an optimal
design and implementation of a complex software system.

\section{How to do abstract thinking?}
Short answer is ``Learn to think concrete first.'' Thinking in high
level requires familiarity with the subject.
Before one can master designing the high-level
architecture, abstract modules, one need experience with instantiation, experience with how each
module would work in the system.
Thinking design Abstract reasoning is based on making up details
beneath the high level concepts.
Value discovery needs the sense of foresight.

\section{Science vs. Engineering}
Science is about beating the wall.
Engineering is about cushioning the wall.

\section[Ideas]{What Ideas do I have}

\begin{itemize}
    \item Know the features of your language and take advantage of them
    \item
\end{itemize}

\section{Coding examples}

\begin{itemize}
    \item Minimum Unique World \url{https://leetcode.com/problems/minimum-unique-word-abbreviation/description/}
    \item
\end{itemize}

\section{Add a chapter for Test coverage}
